% Generated draft: requires mathematical and source review.
\documentclass[11pt]{article}
\usepackage[margin=25mm]{geometry}
\usepackage{fontspec}
\setmainfont{Latin Modern Roman}
\usepackage{amsmath,amssymb,mathtools,unicode-math}
\setmathfont{Latin Modern Math}
\usepackage{enumitem,needspace,xurl,hyperref}
\setlength{\parindent}{0pt}
\setlength{\parskip}{5pt}
\setlength{\emergencystretch}{3em}
\newcommand{\N}{\mathbb{N}}
\newcommand{\Z}{\mathbb{Z}}
\newcommand{\Q}{\mathbb{Q}}
\newcommand{\R}{\mathbb{R}}
\newcommand{\C}{\mathbb{C}}
\newcommand{\F}{\mathbb{F}}
\newcommand{\A}{\mathbb{A}}
\newcommand{\PP}{\mathbb{P}}
\newcommand{\E}{\mathbb{E}}
\newcommand{\eps}{\varepsilon}
\newcommand{\dd}{\,\mathrm{d}}
\begin{document}
\section*{UnsolvedMath ID: 30005242}
\textbf{Source problem code:} catalogue entry
\subsection{Definitions and mathematical statement}
In finite element terminology, the number of degrees of freedom is the dimension \(N=\dim V\) of the finite-dimensional trial space \(V\). The notation \(hp\)-FEM refers to finite element methods involving a mesh-size parameter \(h\) and a polynomial-degree parameter \(p\). These definitions explain the supplied terminology; they do not specify a discretization or a threshold.

The supplied extraction reads ``Is it enough?'' followed by the presentation title \emph{Decompositions of high-frequency Helmholtz solutions and application to the finite element method}. It does not identify what must be sufficient, give a degrees-of-freedom bound, or specify an equation, spatial domain, coefficients, boundary or radiation conditions, admissible data, approximation norm, or quantifiers. Consequently, no precise mathematical assertion or open problem can be recovered from this record.

The supplied statement-verification field says that the official report answers the degrees-of-freedom question and develops its extension, and that the extraction does not retain an open problem. This is a supplied source-verification assessment, not an independently verified conclusion. The supplied literature assessment further says that a degrees-of-freedom threshold is sufficient for \(hp\)-FEM under analytic-obstacle assumptions, while broader variable-coefficient and general-scattering settings remain conditional or incomplete in the cited results. Neither the threshold nor the necessary hypotheses are given, so this assessment cannot be converted here into a quantified theorem or conjecture.
\subsection{Short English statement}
The record concerns degrees-of-freedom requirements for finite element approximation of high-frequency Helmholtz solutions, but the extracted question has lost its mathematical target. Its supplied status is partially solved, while its statement is marked unrecoverable. A precise statement requires recovery and human review of the relevant report passage.
\subsection{Sources}
\begin{itemize}
\item Local registry: \texttt{data/unsolvedmath/problems.json}, with \texttt{catalog/problems.json} designated as the supplementary registry; supplied permanent UnsolvedMath ID \(30005242\), exact source problem code \texttt{catalogue entry}. These are provenance identifiers supplied for this draft; no file inspection is asserted.
\item Supplied report citation: \emph{At the Interface between Semiclassical Analysis and Numerical Analysis of Wave Scattering Problems} (2023), Oberwolfach Reports, DOI \texttt{10.4171/owr/2022/43}, pp.\ 2511--2587. Literature-source label: Oberwolfach Report 43/2022. \url{https://doi.org/10.4171/owr/2022/43}. The extraction includes page number 2551.
\item Supplied presentation attribution: David Lafontaine, \emph{Decompositions of high-frequency Helmholtz solutions and application to the finite element method}, joint work with Jeffrey Galkowski, Euan A. Spence, and Jared Wunsch. The record does not separately identify a problem proposer.
\item Supplied literature reference: Melenk--Sauter, \emph{Mathematics of Computation}. \url{https://doi.org/10.1090/S0025-5718-10-02362-8}. No article title or further bibliographic details are supplied.
\item Literature checked at: \texttt{2026-08-22}. The partially solved label and the assessment concerning analytic obstacles and broader scattering settings are supplied dated curation derived from the supplied OWR curation, not independently verified mathematical facts or a current literature assessment. The record explicitly requests independent verification of time-sensitive status and cited claims.
\end{itemize}
\end{document}
